% ============================================================
%  CHAPTER 21: ELECTROSTATICS
%  The Physics Codex: A Complete University Foundation
%  Korede Samuel Omotosho (Falcon)
%  Aurikrex Learning Systems, 2026
%
%  File: chapters/part4/ch21_electrostatics.tex
% ============================================================

\chapter{Electrostatics}
\label{ch:electrostatics}
\minitoc
\newpage

\chapterquote{Electricity is really just organised
lightning.}{George Carlin}

\begin{objectivebox}
By the end of this chapter, you should be able to:
\begin{enumerate}[noitemsep]
  \item Describe the nature of electric charge and
  the structure of matter that gives rise to it
  \item State and apply Coulomb's Law
  \item Define electric field and electric field
  strength
  \item Sketch electric field patterns for various
  charge configurations
  \item Define electric potential and potential
  difference
  \item Relate electric field strength to potential
  gradient
  \item Describe the behaviour of conductors and
  insulators in electric fields
  \item Explain charging by friction, contact,
  and induction
  \item Describe the gold-leaf electroscope and
  its uses
  \item Apply knowledge of electrostatics to
  practical situations including lightning protection
\end{enumerate}
\end{objectivebox}

% ============================================================
\section{Intuition: What is Electric Charge?}
% ============================================================

Rub a plastic ruler vigorously on your hair, then hold
it near small scraps of paper. The paper scraps leap
toward the ruler and stick to it. Nothing is touching
them. No wind is blowing. Yet they move. Something
invisible is acting across the gap between the ruler
and the paper.

This invisible influence is the \textbf{electric force},
and it arises from a property of matter called
\textbf{electric charge}. Charge is as fundamental as
mass. You cannot explain what charge \textit{is} in
terms of something simpler --- it simply \textit{is},
a fundamental property of certain subatomic particles.
What we can do is describe how it behaves.

Rubbing the ruler on your hair transfers electrons from
your hair to the ruler. The ruler gains excess electrons
and becomes negatively charged. Your hair loses electrons
and becomes positively charged. The attraction between
the charged ruler and the neutral paper is due to the
charged ruler inducing a charge redistribution in the
paper, pulling the opposite charges slightly closer ---
and since opposite charges attract, the net force is
attractive even though the paper was initially neutral.

% ============================================================
\section{Electric Charge}
% ============================================================

\begin{definitionbox}[Electric Charge]
\textbf{Electric charge} $Q$ is a fundamental property
of matter. It exists in two types: \textbf{positive}
and \textbf{negative}.

The SI unit of charge is the \textbf{coulomb} (C).

The charge of one proton: $+e = +1.6\times10^{-19}\ \text{C}$\\
The charge of one electron: $-e = -1.6\times10^{-19}\ \text{C}$\\
Neutrons have zero charge.

Charge is \textbf{quantised}: all observable charges
are integer multiples of the elementary charge $e$:
\[
  Q = ne \quad (n = 0, \pm1, \pm2, \ldots)
\]
\end{definitionbox}

\begin{lawbox}[Conservation of Electric Charge]
The total electric charge in an isolated system is
\textbf{conserved}. Charge cannot be created or
destroyed, only transferred from one object to another.

Like charges \textbf{repel}. Unlike charges
\textbf{attract}.
\end{lawbox}

\begin{diagbox}[Attraction and Repulsion of Charges]
\begin{center}
\begin{tikzpicture}[scale=1.0, font=\small]

  % === UNLIKE CHARGES ATTRACT ===
  \begin{scope}[shift={(0,0)}]
    \node[font=\bfseries] at (2.5,3.5)
      {Unlike charges: ATTRACT};

    % Positive charge (left)
    \fill[codexred] (0.8,1.5) circle (0.5cm);
    \node[white, font=\bfseries] at (0.8,1.5) {$+$};

    % Negative charge (right)
    \fill[codexblue] (4.2,1.5) circle (0.5cm);
    \node[white, font=\bfseries] at (4.2,1.5) {$-$};

    % Force arrows pointing TOWARD each other
    \draw[->, very thick, codexred, line width=2.5pt]
      (1.3,1.5) -- (2.3,1.5);
    \draw[->, very thick, codexblue, line width=2.5pt]
      (3.7,1.5) -- (2.7,1.5);

    % Field lines between charges
    \foreach \y in {1.0,1.5,2.0} {
      \draw[->, codexgold!60, dashed]
        (1.4,\y) -- (3.6,\y);
    }
  \end{scope}

  % === LIKE CHARGES REPEL ===
  \begin{scope}[shift={(7,0)}]
    \node[font=\bfseries] at (2.5,3.5)
      {Like charges: REPEL};

    % Two positive charges
    \fill[codexred] (0.8,1.5) circle (0.5cm);
    \node[white, font=\bfseries] at (0.8,1.5) {$+$};

    \fill[codexred] (4.2,1.5) circle (0.5cm);
    \node[white, font=\bfseries] at (4.2,1.5) {$+$};

    % Force arrows pointing AWAY from each other
    \draw[->, very thick, codexred, line width=2.5pt]
      (1.3,1.5) -- (0.2,1.5);
    \draw[->, very thick, codexred, line width=2.5pt]
      (3.7,1.5) -- (4.8,1.5);

    % Two negative charges
    \fill[codexblue] (0.8,0.0) circle (0.5cm);
    \node[white, font=\bfseries] at (0.8,0.0) {$-$};

    \fill[codexblue] (4.2,0.0) circle (0.5cm);
    \node[white, font=\bfseries] at (4.2,0.0) {$-$};

    \draw[->, very thick, codexblue, line width=2.5pt]
      (1.3,0.0) -- (0.2,0.0);
    \draw[->, very thick, codexblue, line width=2.5pt]
      (3.7,0.0) -- (4.8,0.0);
  \end{scope}

\end{tikzpicture}
\end{center}
\end{diagbox}

% ============================================================
\section{Conductors and Insulators}
% ============================================================

\begin{defbox}[Conductors and Insulators]
\textbf{Conductors} contain large numbers of free
(mobile) electrons that can move throughout the
material. When a charge is placed on a conductor, the
free electrons redistribute until the electric field
inside the conductor is zero. Examples: all metals,
graphite, electrolyte solutions.

\textbf{Insulators} (dielectrics) have electrons
tightly bound to their atoms. Charge placed on an
insulator stays where it is placed. Examples: plastic,
rubber, glass, dry wood, air, ceramics.

\textbf{Semiconductors} have conductivity between
conductors and insulators, controllable by doping.
Examples: silicon, germanium (Chapter 29).
\end{defbox}

\subsection{Methods of Charging}

\begin{diagbox}[Three Methods of Charging]
\begin{center}
\begin{tikzpicture}[scale=0.9, font=\small]

  % === FRICTION ===
  \begin{scope}[shift={(0,5.5)}]
    \node[font=\bfseries] at (3.5,2.0)
      {1. Charging by Friction};
    % Fur (left)
    \fill[brown!30] (0,0) rectangle (2.5,1.2);
    \draw[thick] (0,0) rectangle (2.5,1.2);
    \node at (1.25,0.6) {Fur ($+$)};
    \foreach \x in {0.3,0.7,1.1,1.5,1.9} {
      \node[codexred, font=\bfseries] at (\x,1.0)
        {$+$};
    }
    % Rod (right)
    \fill[codexblue!20] (3.0,0) rectangle (5.5,1.2);
    \draw[thick] (3.0,0) rectangle (5.5,1.2);
    \node at (4.25,0.6) {Rod ($-$)};
    \foreach \x in {3.3,3.7,4.1,4.5,4.9} {
      \node[codexblue, font=\bfseries] at (\x,1.0)
        {$-$};
    }
    % Arrows showing electron transfer
    \draw[->, thick, codexblue]
      (2.8,0.8) -- (3.2,0.8);
    \node[below, font=\tiny] at (3.0,0.2)
      {Electrons transfer from fur to rod\\
      Fur becomes $+$ve, Rod becomes $-$ve};
  \end{scope}

  % === CONTACT ===
  \begin{scope}[shift={(0,2.5)}]
    \node[font=\bfseries] at (3.5,2.0)
      {2. Charging by Contact};
    % Charged sphere
    \fill[codexred!40] (0.8,0.6) circle (0.7);
    \draw[thick] (0.8,0.6) circle (0.7);
    \foreach \angle in {0,45,90,135,180,225,270,315} {
      \node[codexred, font=\tiny]
        at ({0.8+0.5*cos(\angle)},{0.6+0.5*sin(\angle)})
        {$+$};
    }
    \node[below] at (0.8,-0.2) {Charged ($+$)};

    % Arrow
    \draw[->, thick, codexred] (1.6,0.6) -- (2.5,0.6);
    \node[above, font=\tiny] at (2.0,0.7) {touch};

    % Neutral sphere
    \fill[gray!20] (3.3,0.6) circle (0.7);
    \draw[thick] (3.3,0.6) circle (0.7);
    \node[below] at (3.3,-0.2) {Neutral};

    % Arrow
    \draw[->, thick, codexred] (4.4,0.6) -- (5.3,0.6);
    \node[above, font=\tiny] at (4.8,0.7) {separate};

    % Both charged
    \fill[codexred!30] (5.8,0.6) circle (0.7);
    \draw[thick] (5.8,0.6) circle (0.7);
    \foreach \angle in {0,90,180,270} {
      \node[codexred, font=\tiny]
        at ({5.8+0.4*cos(\angle)},{0.6+0.4*sin(\angle)})
        {$+$};
    }
    \node[below] at (5.8,-0.2) {Also $+$ve};
  \end{scope}

  % === INDUCTION ===
  \begin{scope}[shift={(0,0)}]
    \node[font=\bfseries] at (3.5,2.0)
      {3. Charging by Induction (no contact)};
    % Charged rod nearby
    \fill[codexred!30] (-0.5,0.3) rectangle (0.0,1.3);
    \foreach \y in {0.5,0.8,1.1} {
      \node[codexred, font=\tiny] at (-0.25,\y) {$+$};
    }
    \node[left, font=\tiny] at (-0.6,0.8) {Charged\\rod};

    % Neutral conductor with induced charges
    \fill[gray!20] (0.8,0) rectangle (4.2,1.6);
    \draw[thick] (0.8,0) rectangle (4.2,1.6);
    % Negative induced charges (near rod)
    \foreach \y in {0.3,0.6,0.9,1.2} {
      \node[codexblue, font=\bfseries] at (1.1,\y) {$-$};
    }
    % Positive induced charges (far end)
    \foreach \y in {0.3,0.6,0.9,1.2} {
      \node[codexred, font=\bfseries] at (3.9,\y) {$+$};
    }
    \node[below, font=\tiny] at (2.5,-0.2)
      {Conductor: charge redistributes by induction};

    % Earth the far end, remove earth, remove rod
    \draw[->, thick, codexblue] (4.2,0.8) -- (5.0,0.8)
      node[right, font=\tiny]{Earth\\(remove $+$)};
    \fill[gray!20] (5.8,0) rectangle (7.5,1.6);
    \draw[thick] (5.8,0) rectangle (7.5,1.6);
    \foreach \y in {0.3,0.6,0.9,1.2} {
      \node[codexblue, font=\bfseries] at (6.0,\y)
        {$-$};
    }
    \node[below, font=\tiny] at (6.65,-0.2)
      {Conductor now $-$ve\\(no contact with rod)};
  \end{scope}

\end{tikzpicture}
\end{center}
\end{diagbox}

% ============================================================
\section{Coulomb's Law}
% ============================================================

\begin{historybox}[Charles-Augustin de Coulomb, 1785]
Charles Coulomb (1736--1806) used a torsion balance
(similar to the one later used by Cavendish to measure
$G$) to measure the force between charged spheres with
extraordinary precision. He established that the
electrostatic force follows an inverse square law ---
exactly parallel to Newton's gravitational law --- and
that it depends on the product of the charges. His
careful measurements in 1785 remain one of the most
elegant experimental confirmations of a fundamental
physical law.
\end{historybox}

\begin{lawbox}[Coulomb's Law]
The electrostatic force between two point charges
$Q_1$ and $Q_2$ separated by a distance $r$ in
vacuum is:
\[
  F = \frac{kQ_1Q_2}{r^2}
  = \frac{Q_1Q_2}{4\pi\varepsilon_0 r^2}
\]
where:
\begin{itemize}[noitemsep]
  \item $k = 8.99\times10^9\ \text{N\,m}^2\text{C}^{-2}$
  is Coulomb's constant
  \item $\varepsilon_0 = 8.85\times10^{-12}\ \text{C}^2\text{N}^{-1}\text{m}^{-2}$
  is the permittivity of free space
  \item $F > 0$: repulsive force (like charges)
  \item $F < 0$: attractive force (unlike charges)
\end{itemize}
The force is along the line joining the two charges
(central force).
\end{lawbox}

\begin{keybox}[Comparison: Coulomb's Law vs Newton's Gravity]
\begin{center}
\begin{tabular}{lll}
\toprule
\textbf{Feature} & \textbf{Coulomb} & \textbf{Newton}\\
\midrule
Formula & $F = kQ_1Q_2/r^2$ & $F = Gm_1m_2/r^2$\\
Distance dependence & $1/r^2$ & $1/r^2$\\
Both attractive? & No (can repel) & Yes (always)\\
Constant & $k = 8.99\times10^9$ &
  $G = 6.67\times10^{-11}$\\
Magnitude ratio & $F_E/F_G \approx 10^{39}$ (for electrons) & \\
\bottomrule
\end{tabular}
\end{center}
\vspace{0.3cm}
The electrostatic force between two electrons is
$\sim10^{39}$ times stronger than the gravitational
force between them. Gravity is the dominant force
on cosmic scales only because matter is electrically
neutral in bulk.
\end{keybox}

\begin{examplebox}[21.1]
\stepcounter{workedexample}
Two point charges, $Q_1 = +3.0\ \mu\text{C}$ and
$Q_2 = -5.0\ \mu\text{C}$, are placed $0.20\ \text{m}$
apart.\\
(a) Calculate the force between them.\\
(b) Is the force attractive or repulsive?\\
(c) How does the force change if the distance is
tripled?\\
($k = 9.0\times10^9\ \text{N\,m}^2\text{C}^{-2}$)
\end{examplebox}

\begin{solutionbox}
\textbf{(a)} Force:
\[
  F = \frac{kQ_1Q_2}{r^2}
  = \frac{9.0\times10^9 \times 3.0\times10^{-6}
    \times 5.0\times10^{-6}}{(0.20)^2}
  = \frac{9.0\times10^9 \times 1.5\times10^{-11}}
    {0.04}
\]
\[
  = \frac{0.135}{0.04} = 3.375\ \text{N}
  \approx 3.4\ \text{N}
\]

\textbf{(b)} The charges are opposite ($+$ and $-$)
so the force is \textbf{attractive}. The negative
sign in the product $Q_1Q_2$ indicates attraction.

\textbf{(c)} When $r$ is tripled ($r' = 3r$):
\[
  F' = \frac{kQ_1Q_2}{(3r)^2}
  = \frac{F}{9} = \frac{3.4}{9} \approx 0.38\ \text{N}
\]
Tripling the distance reduces the force to
$\frac{1}{9}$ of its original value (inverse square
law).
\end{solutionbox}

% ============================================================
\section{Electric Field}
% ============================================================

\begin{definitionbox}[Electric Field Strength]
An \textbf{electric field} exists in the region
surrounding a charged object. It is a field of force.

The \textbf{electric field strength} $\vect{E}$ at a
point is the force per unit positive charge placed
at that point:
\[
  \vect{E} = \frac{\vect{F}}{q}
  \implies F = qE
\]
SI unit: N\,C$^{-1}$ or V\,m$^{-1}$ (both equivalent).

Direction of $\vect{E}$: same as the force on a
\textit{positive} test charge.
\end{definitionbox}

\begin{lawbox}[Electric Field Due to a Point Charge]
The electric field strength at distance $r$ from a
point charge $Q$ in vacuum:
\[
  E = \frac{kQ}{r^2} = \frac{Q}{4\pi\varepsilon_0 r^2}
\]
Direction: radially outward for positive $Q$;
radially inward for negative $Q$.
\end{lawbox}

\begin{examplebox}[21.2]
\stepcounter{workedexample}
A proton ($Q = 1.6\times10^{-19}\ \text{C}$) is
placed in a uniform electric field of
$E = 5.0\times10^4\ \text{N\,C}^{-1}$.\\
(a) Find the force on the proton.\\
(b) Find the acceleration of the proton.
($m_p = 1.67\times10^{-27}\ \text{kg}$)\\
(c) Find the field strength at a distance of
$0.10\ \text{m}$ from a $+2.0\ \mu\text{C}$
point charge in vacuum.
\end{examplebox}

\begin{solutionbox}
\textbf{(a)} Force on proton:
\[
  F = qE = 1.6\times10^{-19}
  \times 5.0\times10^4
  = 8.0\times10^{-15}\ \text{N}
\]

\textbf{(b)} Acceleration:
\[
  a = \frac{F}{m} = \frac{8.0\times10^{-15}}
  {1.67\times10^{-27}}
  = 4.79\times10^{12}\ \text{m\,s}^{-2}
\]
About $5\times10^{11}\ g$ --- an enormous acceleration.
Electric forces on subatomic particles are vastly
more significant than gravity.

\textbf{(c)}
\[
  E = \frac{kQ}{r^2}
  = \frac{9.0\times10^9 \times 2.0\times10^{-6}}
    {(0.10)^2}
  = \frac{1.8\times10^4}{0.01}
  = 1.8\times10^6\ \text{N\,C}^{-1}
\]
\end{solutionbox}

% ============================================================
\section{Electric Field Lines}
% ============================================================

Electric fields are visualised using \textbf{field
lines} (lines of force). Rules for drawing field lines:
\begin{enumerate}[noitemsep]
  \item Field lines start on positive charges and
  end on negative charges (or extend to/from infinity).
  \item The direction of a field line at any point is
  the direction of $\vect{E}$ at that point.
  \item The density of field lines (lines per unit
  area) represents the magnitude of $\vect{E}$.
  \item Field lines never cross (there is only one
  direction of $\vect{E}$ at any point).
  \item Field lines are perpendicular to conducting
  surfaces.
\end{enumerate}

\begin{diagbox}[Electric Field Patterns for Common Charge Configurations]
\begin{center}
\begin{tikzpicture}[scale=0.9, font=\small]

  % === SINGLE POSITIVE CHARGE ===
  \begin{scope}[shift={(0,6)}]
    \node[font=\bfseries] at (0,2.2) {$+Q$};
    \fill[codexred] (0,0) circle (0.35cm);
    \node[white, font=\bfseries] at (0,0) {$+$};
    % Radial field lines out
    \foreach \angle in {0,45,90,135,180,225,270,315} {
      \draw[->, thick, codexred]
        ({0.4*cos(\angle)},{0.4*sin(\angle)})
        -- ({1.8*cos(\angle)},{1.8*sin(\angle)});
    }
  \end{scope}

  % === SINGLE NEGATIVE CHARGE ===
  \begin{scope}[shift={(5,6)}]
    \node[font=\bfseries] at (0,2.2) {$-Q$};
    \fill[codexblue] (0,0) circle (0.35cm);
    \node[white, font=\bfseries] at (0,0) {$-$};
    % Radial field lines in
    \foreach \angle in {0,45,90,135,180,225,270,315} {
      \draw[->, thick, codexblue]
        ({1.8*cos(\angle)},{1.8*sin(\angle)})
        -- ({0.4*cos(\angle)},{0.4*sin(\angle)});
    }
  \end{scope}

  % === TWO OPPOSITE CHARGES (dipole) ===
  \begin{scope}[shift={(0,1)}]
    \node[font=\bfseries] at (2.5,3.5)
      {$+Q$ and $-Q$ (dipole field)};
    \fill[codexred] (0.5,1.5) circle (0.35cm);
    \node[white, font=\bfseries] at (0.5,1.5) {$+$};
    \fill[codexblue] (4.5,1.5) circle (0.35cm);
    \node[white, font=\bfseries] at (4.5,1.5) {$-$};

    % Field lines from + to -
    % Straight horizontal line
    \draw[->, thick, codexgold]
      (0.85,1.5) -- (4.15,1.5);
    % Curved lines
    \draw[->, thick, codexgold]
      (0.5,1.85) .. controls (1.5,3.0) and (3.5,3.0)
      .. (4.5,1.85);
    \draw[->, thick, codexgold]
      (0.5,1.15) .. controls (1.5,0.0) and (3.5,0.0)
      .. (4.5,1.15);
    \draw[->, thick, codexgold!70]
      (0.85,1.9) .. controls (2.0,3.8) and (3.0,3.8)
      .. (4.15,1.9);
    \draw[->, thick, codexgold!70]
      (0.85,1.1) .. controls (2.0,-0.8) and (3.0,-0.8)
      .. (4.15,1.1);
  \end{scope}

  % === UNIFORM FIELD (parallel plates) ===
  \begin{scope}[shift={(7,0)}]
    \node[font=\bfseries] at (2.5,4.5)
      {Uniform field (parallel plates)};
    % Plates
    \fill[codexred!40] (0,3.5) rectangle (5,4.0);
    \fill[codexblue!40] (0,0) rectangle (5,0.5);
    \draw[thick] (0,3.5) -- (5,3.5);
    \draw[thick] (0,0.5) -- (5,0.5);
    % $+$ on top plate
    \foreach \x in {0.5,1.5,2.5,3.5,4.5} {
      \node[codexred, font=\bfseries] at (\x,3.75)
        {$+$};
    }
    % $-$ on bottom plate
    \foreach \x in {0.5,1.5,2.5,3.5,4.5} {
      \node[codexblue, font=\bfseries] at (\x,0.25)
        {$-$};
    }
    % Uniform downward field lines
    \foreach \x in {0.6,1.3,2.0,2.7,3.4,4.1,4.8} {
      \draw[->, thick, codexgold]
        (\x,3.5) -- (\x,0.5);
    }
    \node[right, font=\tiny, codexgold] at (5.1,2.0)
      {Uniform $E$\\(parallel lines,\\equal spacing)};
    % E field label
    \node[codexgold, font=\small] at (2.5,2.0) {$\vect{E}$};
    \draw[->, very thick, codexgold]
      (2.5,2.5) -- (2.5,1.5);
  \end{scope}

\end{tikzpicture}
\end{center}
\end{diagbox}

% ============================================================
\section{Uniform Electric Field}
% ============================================================

Between two parallel conducting plates separated by
distance $d$ and with potential difference $V$ between
them, the electric field is \textbf{uniform}
(constant magnitude and direction):

\begin{lawbox}[Uniform Electric Field Between Parallel Plates]
\[
  E = \frac{V}{d}
\]
where $V$ is the potential difference between the
plates and $d$ is the separation.

The force on a charge $q$ in this field:
\[
  F = qE = \frac{qV}{d}
\]
\end{lawbox}

% ============================================================
\section{Electric Potential}
% ============================================================

\begin{definitionbox}[Electric Potential and Potential
Difference]
The \textbf{electric potential} $V$ at a point is the
work done per unit positive charge in bringing a small
positive test charge from infinity (where $V = 0$) to
that point:
\[
  V = \frac{W}{q}
\]
SI unit: volt (V) where $1\ \text{V}
= 1\ \text{J\,C}^{-1}$.

The \textbf{potential difference} (p.d.) between two
points A and B is:
\[
  V_{AB} = V_A - V_B
  = \frac{W_{AB}}{q}
\]
where $W_{AB}$ is the work done moving charge $q$
from B to A.

For a point charge $Q$ in vacuum:
\[
  V = \frac{kQ}{r} = \frac{Q}{4\pi\varepsilon_0 r}
\]
\end{definitionbox}

\begin{keybox}[Relationship Between $E$ and $V$]
The electric field is related to the potential by:
\[
  E = -\frac{dV}{dr}
  \approx -\frac{\Delta V}{\Delta r}
\]
$E$ is the negative of the potential gradient.
Field lines point from high potential to low potential
(downhill in potential, just as water flows downhill
in gravitational potential).

For a uniform field: $E = V/d$.
\end{keybox}

\begin{definitionbox}[Equipotential Surfaces]
\textbf{Equipotential surfaces} are surfaces on
which the electric potential is the same at every
point. No work is done moving a charge along an
equipotential surface.

Equipotential surfaces are always \textbf{perpendicular}
to electric field lines.

For a point charge: equipotentials are spheres
centred on the charge.
For a uniform field: equipotentials are planes
perpendicular to $\vect{E}$.
\end{definitionbox}

\begin{diagbox}[Field Lines and Equipotentials for a Point Charge]
\begin{center}
\begin{tikzpicture}[scale=1.0, font=\small]

  % Charge at centre
  \fill[codexred] (0,0) circle (0.3cm);
  \node[white, font=\bfseries\small] at (0,0) {$+$};

  % Radial field lines (8 directions)
  \foreach \angle in {0,45,90,135,180,225,270,315} {
    \draw[->, very thick, codexred!70]
      ({0.35*cos(\angle)},{0.35*sin(\angle)})
      -- ({2.5*cos(\angle)},{2.5*sin(\angle)});
  }

  % Equipotential circles
  \foreach \r/\label in {
    0.7/{$V_1$}, 1.2/{$V_2$}, 1.8/{$V_3$}} {
    \draw[codexblue, thick, dashed] (0,0) circle (\r);
    \node[codexblue, font=\tiny] at (\r+0.2,0.15)
      {\label};
  }

  % Annotation
  \node[font=\small, align=center] at (0,-3.2)
    {Radial field lines $\perp$ to circular\\
    equipotentials. $V_1 > V_2 > V_3$\\
    (potential falls with distance from $+$ charge)};

  % Key
  \draw[codexred!70, thick, ->] (3.0,1.0) -- (3.7,1.0);
  \node[right, codexred!70, font=\tiny] at (3.8,1.0)
    {Field lines};
  \draw[codexblue, thick, dashed]
    (3.0,0.4) arc(0:30:0.5);
  \node[right, codexblue, font=\tiny] at (3.6,0.4)
    {Equipotentials};

\end{tikzpicture}
\end{center}
\end{diagbox}

\begin{examplebox}[21.3]
\stepcounter{workedexample}
Two parallel plates are separated by $4.0\ \text{cm}$
with a potential difference of $800\ \text{V}$.
A positive charge of $3.0\times10^{-6}\ \text{C}$
is placed between the plates.\\
(a) Find the electric field strength.\\
(b) Find the force on the charge.\\
(c) Find the work done moving the charge from the
negative plate to the positive plate.\\
(d) Find the potential at a point halfway between
the plates, taking the negative plate as zero
potential.
\end{examplebox}

\begin{solutionbox}
\textbf{(a)} Electric field:
\[
  E = \frac{V}{d} = \frac{800}{0.040}
  = 20\,000\ \text{V\,m}^{-1}
  = 2.0\times10^4\ \text{N\,C}^{-1}
\]

\textbf{(b)} Force:
\[
  F = qE = 3.0\times10^{-6} \times 2.0\times10^4
  = 0.060\ \text{N}
\]
Direction: toward the negative plate (force on
positive charge is in direction of $\vect{E}$,
from high to low potential --- the force on a
positive charge is from the positive plate toward
the negative plate).

Wait --- let us clarify: the field points from the
$+$ plate to the $-$ plate, and the force on a
positive charge is in the direction of $\vect{E}$,
i.e.\ toward the negative plate.

\textbf{(c)} Work done (against electric force,
moving from $-$ to $+$ plate):
\[
  W = qV = 3.0\times10^{-6} \times 800
  = 2.4\times10^{-3}\ \text{J}
\]

\textbf{(d)} Halfway ($d/2 = 0.020\ \text{m}$
from negative plate):
\[
  V_{half} = E \times \frac{d}{2}
  = 20\,000 \times 0.020 = 400\ \text{V}
\]
The potential increases linearly from 0 to 800 V
across the gap.
\end{solutionbox}

% ============================================================
\section{The Gold-Leaf Electroscope}
% ============================================================

\begin{diagbox}[The Gold-Leaf Electroscope]
\begin{center}
\begin{tikzpicture}[scale=1.0, font=\small]

  % Glass jar
  \draw[thick, gray!60] (1.0,0) -- (0.5,4.0);
  \draw[thick, gray!60] (5.0,0) -- (5.5,4.0);
  \draw[thick, gray!60] (1.0,0) -- (5.0,0);
  \draw[thick, gray!60] (0.5,4.0) -- (5.5,4.0);
  \node[gray, font=\itshape, font=\small] at (3,0.5)
    {Glass jar (insulating)};

  % Metal rod
  \fill[gray!50] (2.8,1.5) rectangle (3.2,6.5);
  \draw[thick] (2.8,1.5) rectangle (3.2,6.5);
  \node[right, font=\tiny] at (3.3,5.5) {Metal rod};

  % Metal cap (top)
  \fill[gray!50] (2.0,6.5) rectangle (4.0,7.0);
  \draw[thick] (2.0,6.5) rectangle (4.0,7.0);
  \node[above, font=\small] at (3.0,7.0)
    {Metal cap\\(collecting plate)};

  % Two gold leaves (uncharged: hanging vertical)
  \draw[very thick, codexgold]
    (3.0,1.5) -- (2.3,0.2);
  \draw[very thick, codexgold]
    (3.0,1.5) -- (3.7,0.2);
  \node[below, codexgold] at (2.0,0.1)
    {Gold leaf};
  \node[below, codexgold] at (4.0,0.1)
    {Gold leaf};

  % Charged version (leaves spread)
  \begin{scope}[shift={(7,0)}]
    \draw[thick, gray!60] (1.0,0) -- (0.5,4.0);
    \draw[thick, gray!60] (5.0,0) -- (5.5,4.0);
    \draw[thick, gray!60] (1.0,0) -- (5.0,0);
    \draw[thick, gray!60] (0.5,4.0) -- (5.5,4.0);
    \fill[gray!50] (2.8,1.5) rectangle (3.2,6.5);
    \draw[thick] (2.8,1.5) rectangle (3.2,6.5);
    \fill[gray!50] (2.0,6.5) rectangle (4.0,7.0);
    \draw[thick] (2.0,6.5) rectangle (4.0,7.0);
    \node[above, font=\small] at (3.0,7.0)
      {Charged metal cap};
    % Charge signs on cap
    \foreach \x in {2.2,2.7,3.2,3.7} {
      \node[codexred, font=\bfseries] at (\x,6.75)
        {$+$};
    }
    % Charge signs on rod
    \foreach \y in {2.0,2.8,3.6,4.4,5.2,6.0} {
      \node[codexred, font=\bfseries] at (3.0,\y)
        {$+$};
    }
    % Leaves spread apart
    \draw[very thick, codexgold]
      (3.0,1.5) -- (1.8,-0.2);
    \draw[very thick, codexgold]
      (3.0,1.5) -- (4.2,-0.2);
    % Charge on leaves
    \node[codexred, font=\bfseries] at (2.1,0.6)
      {$+$};
    \node[codexred, font=\bfseries] at (3.9,0.6)
      {$+$};
    \node[font=\small, codexred, align=center]
      at (3.0,-0.7)
      {Leaves repel:\\same sign charge};
  \end{scope}

  % Labels
  \node[font=\small, align=center] at (3.0,-1.5)
    {Uncharged: leaves hang\\vertically (no force)};

\end{tikzpicture}
\end{center}
\end{diagbox}

The gold-leaf electroscope detects the presence of
charge. When a charged object is brought near or
touches the cap, the leaves diverge because they
acquire the same sign of charge and repel each other.
The greater the charge, the larger the divergence.

\textbf{Uses of the electroscope:}
\begin{enumerate}[noitemsep]
  \item Detecting the presence and sign of charge
  \item Testing whether a body is a conductor or
  insulator
  \item Detecting ionising radiation (charged
  particles discharge the electroscope)
\end{enumerate}

% ============================================================
\section{Properties of Conductors in Electrostatic
Equilibrium}
% ============================================================

\begin{lawbox}[Properties of Charged Conductors]
For a conductor in electrostatic equilibrium (charges
not moving):
\begin{enumerate}[noitemsep]
  \item The electric field \textbf{inside} the
  conductor is \textbf{zero}.
  \item Any excess charge resides entirely on the
  \textbf{outer surface} of the conductor.
  \item The electric field at the surface is
  perpendicular to the surface.
  \item Charge concentrates on sharp points and
  edges (where curvature is greatest).
  \item The entire conductor is at the same
  potential (equipotential body).
\end{enumerate}
\end{lawbox}

\begin{realworldbox}[Lightning Conductors]
Lightning is a massive electrostatic discharge
between charged clouds and the Earth. A lightning
conductor (lightning rod) on a tall building works
by exploiting two properties: (1) the sharp metal
spike at the top concentrates electric field lines,
ionising the air and providing a conductive path;
(2) the copper strip connecting the spike to earth
provides a low-resistance path for the current of
billions of amperes to flow safely to ground,
bypassing the building structure.

The \textbf{National Theatre} in Lagos, the
\textbf{Transcorp Hilton} in Abuja, and all tall
structures in Nigeria are required by building codes
to have lightning protection systems. Benjamin
Franklin invented the lightning rod in 1752, the
same year he famously flew a kite in a thunderstorm
to demonstrate that lightning was electrical. He
survived. Several who attempted to repeat his
experiment did not.
\end{realworldbox}

% ============================================================
\section{Gauss's Law (Qualitative)}
% ============================================================

\begin{lawbox}[Gauss's Law (Statement)]
The total electric flux $\Phi_E$ through any closed
surface is proportional to the total charge enclosed
within that surface:
\[
  \Phi_E = \oint \vect{E} \cdot d\vect{A}
  = \frac{Q_{enclosed}}{\varepsilon_0}
\]
This powerful law allows us to calculate electric
fields for symmetric charge distributions very
efficiently. Key results:

\textbf{Outside a spherical conductor (radius $R$):}
$E = kQ/r^2$ (same as point charge)

\textbf{Inside a spherical conductor:} $E = 0$

\textbf{Uniform infinite plate (surface charge
density $\sigma$):} $E = \sigma/(2\varepsilon_0)$

\textbf{Between parallel plates:}
$E = \sigma/\varepsilon_0 = V/d$
\end{lawbox}

% ============================================================
\section{Chapter Summary}
% ============================================================

\begin{summarybox}
\textbf{Charge:} positive or negative, quantised
in units of $e = 1.6\times10^{-19}\ \text{C}$.
Conserved. Like charges repel; unlike attract.

\textbf{Coulomb's Law:}
$F = kQ_1Q_2/r^2$; $k = 9.0\times10^9\ \text{N\,m}^2\text{C}^{-2}$

\textbf{Electric field:} $E = F/q$ (N\,C$^{-1}$).
Point charge: $E = kQ/r^2$.
Uniform field: $E = V/d$.

\textbf{Electric potential:} $V = W/q$ (volts).
Point charge: $V = kQ/r$.
Potential decreases in direction of $\vect{E}$.
$E = -\Delta V/\Delta r$.

\textbf{Equipotentials:} surfaces of equal $V$,
always $\perp$ to field lines.

\textbf{Conductors:} $E = 0$ inside; all charge on
surface; surface is equipotential; charge
concentrates at sharp points.

\textbf{Charging methods:} friction (electron
transfer); contact (charge sharing); induction
(no contact needed).

\textbf{Gauss's Law:}
$\Phi = Q_{enc}/\varepsilon_0$.
Explains why $E = 0$ inside a conductor.
\end{summarybox}

% ============================================================
\newpage
\begin{exercisebox}[21]

\subsection*{Section A: Multiple Choice}

\textbf{A1.} Two charges $+2Q$ and $+3Q$ are placed
$d$ apart. The force between them is $F$. If the
distance is halved and each charge is doubled, the
new force is:

\mcoption{A}{$F$}
\mcoption{B}{$4F$}
\mcoption{C}{$16F$}
\mcoption{D}{$24F$}
\ansref{Ch21-A1}

\vspace{0.4cm}

\textbf{A2.} The electric field inside a uniformly
charged hollow metal sphere is:

\mcoption{A}{Maximum at the centre}
\mcoption{B}{Zero throughout}
\mcoption{C}{Equal to $kQ/r^2$ at all interior points}
\mcoption{D}{Directed radially inward}
\ansref{Ch21-A2}

\vspace{0.4cm}

\textbf{A3.} A charge $q$ moves from a point of
potential $V_1$ to a point of potential $V_2$. The
work done by the electric field on the charge is:

\mcoption{A}{$q(V_2 - V_1)$}
\mcoption{B}{$q(V_1 - V_2)$}
\mcoption{C}{$q(V_1 + V_2)$}
\mcoption{D}{$qV_1V_2$}
\ansref{Ch21-A3}

\vspace{0.4cm}

\textbf{A4.} Two parallel plates are $5.0\ \text{mm}$
apart with a potential difference of $250\ \text{V}$.
The electric field between them is:

\mcoption{A}{$50\ \text{V\,m}^{-1}$}
\mcoption{B}{$500\ \text{V\,m}^{-1}$}
\mcoption{C}{$5\times10^4\ \text{V\,m}^{-1}$}
\mcoption{D}{$5\times10^5\ \text{V\,m}^{-1}$}
\ansref{Ch21-A4}

\vspace{0.4cm}

\textbf{A5.} The electric potential due to a point
charge $Q$ at distance $r$ is proportional to:

\mcoption{A}{$1/r^2$}
\mcoption{B}{$1/r$}
\mcoption{C}{$r$}
\mcoption{D}{$r^2$}
\ansref{Ch21-A5}

\vspace{0.4cm}

\textbf{A6.} A gold-leaf electroscope is negatively
charged. When a positively charged rod is brought
near (but not touching) the cap, the leaves:

\mcoption{A}{Diverge further}
\mcoption{B}{Come together}
\mcoption{C}{Remain unchanged}
\mcoption{D}{The cap becomes positive}
\ansref{Ch21-A6}

\vspace{0.4cm}

\textbf{A7.} Charging by induction produces a
charge that is:

\mcoption{A}{The same sign as the inducing charge}
\mcoption{B}{Opposite sign to the inducing charge}
\mcoption{C}{Zero charge}
\mcoption{D}{Depends on whether the object is a
conductor}
\ansref{Ch21-A7}

\vspace{0.4cm}

\textbf{A8.} The unit of electric field strength
can also be written as:

\mcoption{A}{J\,C$^{-1}$}
\mcoption{B}{V\,m$^{-1}$}
\mcoption{C}{C\,N$^{-1}$}
\mcoption{D}{J\,m$^{-1}$}
\ansref{Ch21-A8}

\vspace{0.4cm}

\textbf{A9.} A proton and an electron are placed
$1\ \text{nm}$ apart. The ratio of the electric
force to the gravitational force between them is
approximately:

\mcoption{A}{$10^{20}$}
\mcoption{B}{$10^{30}$}
\mcoption{C}{$10^{39}$}
\mcoption{D}{$10^{10}$}
\ansref{Ch21-A9}

\vspace{0.4cm}

\textbf{A10.} Equipotential surfaces and electric
field lines are always:

\mcoption{A}{Parallel to each other}
\mcoption{B}{Perpendicular to each other}
\mcoption{C}{Identical}
\mcoption{D}{Parallel only for a point charge}
\ansref{Ch21-A10}

\vspace{0.6cm}

\subsection*{Section B: Theory and Calculations}

\textbf{B1.} Three point charges are arranged as
follows: $Q_1 = +4.0\ \mu\text{C}$ at the origin,
$Q_2 = -2.0\ \mu\text{C}$ at $x = 0.30\ \text{m}$,
$Q_3 = +3.0\ \mu\text{C}$ at $x = 0.60\ \text{m}$.
All on the x-axis.
($k = 9.0\times10^9\ \text{N\,m}^2\text{C}^{-2}$)\\[4pt]
(a) Calculate the force on $Q_2$ due to $Q_1$.\\
(b) Calculate the force on $Q_2$ due to $Q_3$.\\
(c) Find the net force on $Q_2$ (magnitude and
direction).\\
(d) At what point on the x-axis (between $Q_1$ and
$Q_3$) is the electric field zero?
\hfill\ansref{Ch21-B1}

\vspace{0.4cm}

\textbf{B2.} A parallel plate capacitor has plates
of area $0.050\ \text{m}^2$ separated by
$2.0\ \text{mm}$. A potential difference of
$400\ \text{V}$ is applied.\\[4pt]
(a) Calculate the electric field strength between
the plates.\\
(b) Calculate the surface charge density $\sigma$
on the plates (use $E = \sigma/\varepsilon_0$).\\
(c) Calculate the charge on each plate.\\
(d) An electron (mass $9.11\times10^{-31}\ \text{kg}$,
charge $-1.6\times10^{-19}\ \text{C}$) is released
from rest near the negative plate. Find its
acceleration and the speed with which it hits the
positive plate.
\hfill\ansref{Ch21-B2}

\vspace{0.4cm}

\textbf{B3.} A metallic sphere of radius $10\ \text{cm}$
carries a charge of $+5.0\ \mu\text{C}$ in free space.
($k = 9.0\times10^9\ \text{N\,m}^2\text{C}^{-2}$,
$\varepsilon_0 = 8.85\times10^{-12}\ \text{C}^2\text{N}^{-1}\text{m}^{-2}$)\\[4pt]
(a) Calculate the electric field strength at the
surface of the sphere.\\
(b) Calculate the electric field at $0.30\ \text{m}$
from the centre of the sphere.\\
(c) Calculate the electric potential at the surface.\\
(d) Calculate the work done bringing a $+2.0\ \mu\text{C}$
charge from infinity to the surface of the sphere.\\
(e) Explain why the electric field inside the sphere
is zero, using Gauss's Law reasoning.
\hfill\ansref{Ch21-B3}

\vspace{0.4cm}

\textbf{B4.} A charged oil drop of mass
$3.27\times10^{-15}\ \text{kg}$ is held stationary
between two horizontal parallel plates separated
by $6.0\ \text{mm}$. The lower plate is negative.
The potential difference is $490\ \text{V}$.
($g = 10\ \text{m\,s}^{-2}$,
$e = 1.6\times10^{-19}\ \text{C}$)\\[4pt]
(a) Draw a free-body diagram of the stationary
drop showing all forces.\\
(b) For the drop to be stationary, what must be
the relationship between the electric force and
the gravitational force?\\
(c) Calculate the charge on the drop.\\
(d) How many elementary charges does this represent?\\
(e) The potential difference is doubled. Describe
the subsequent motion of the drop.
\hfill\ansref{Ch21-B4}

\vspace{0.4cm}

\textbf{B5.} (a) Explain the principle of operation
of a lightning conductor, referring to:
(i) the distribution of charge on conductors with
sharp points; (ii) the ionisation of air; (iii) the
path of the lightning discharge.\\[4pt]
(b) A lightning strike delivers a charge of
$5.0\ \text{C}$ in $0.001\ \text{s}$.
Calculate the average current during the strike.\\[4pt]
(c) If the potential difference between the cloud
and ground just before the strike is
$3.0\times10^8\ \text{V}$, calculate the energy
released in the lightning strike.\\[4pt]
(d) Express this energy in kWh. How long could this
energy power a $60\ \text{W}$ light bulb?\\[4pt]
(e) Despite the enormous energy in a lightning
strike, it is not practical to harvest lightning
for electricity generation in Nigeria. Suggest
two physical reasons why.
\hfill\ansref{Ch21-B5}

\end{exercisebox}